\documentclass[twocolumn]{notes}
\usepackage{notes}

\begin{document}
\title{Discrete-state kinetics and Markov models}
\date{\today} % \formatdate{10}{11}{2022}
\maketitle
\tableofcontents
\newpage

\section{Discretized Kinetics: Simple examples and transition matrix basics}

\href{https://www.physicallensonthecell.org/discrete-state-kinetics-and-markov-models}{\color{darkred}Zuckermans physical lens article.}

\subsection{Overview}
\begin{itemizec}
    \item life: non-equilibrium = active transport, etc.
    \item physics: non-equilibrium steady state
    \item continuum processes: Fokker-Planck/Smoluchowski
    \item these equations hard to solve!
    \item go to discrete system
    \item simplest discrete system = two state
\end{itemizec}
\subsection{Continuous vs. discretized time}
\begin{itemizec}
    \item continuous time: standard ODEs of chemical kinetics
    \item discrete time: simulation, timesteps
\end{itemizec}

\subsection{Continuous time two state system}
\begin{itemizec}
    \item mass action equations
    \begin{align*}
    \myder{p_0}{t} &= -k_+p_0 + k_-p_1 \\
    \myder{p_1}{t} &= k_+p_0 - k_-p_1
    \end{align*}
    equivalent to
    \[
        \myder{}{t}p(t) = \begin{bmatrix}
            k_+ & k_- \\ k_+ & -k_-
            \end{bmatrix} p(t)
    \]
    \item assumes rate constants are independent of concentration (ignore aggregation, crowding effects)
    \item summing mass action equations gives
    \[
        \myder{}{t}[p_0 + p_1] = 0
    \]
    which is the mass conservation condition
    \[
        p_0 + p_1 = 1
    \]
    \item therefore, solving one equation is sufficient in this case.
    \begin{align*}
        \myder{p_0}{t} = -k_+p_0 + k_-p_1 \\
        \myder{p_0}{t} + \underbrace{[k_+ + k_-]}_{k ~\equiv~ \tau_\mathrm{rel}^{-1}}p_0 - k_- = 0
    \end{align*}
    ane we have a first order \emph{linear} ODE
    \item use the method of \emph{integrating factors} to solve equations of the form
    \[
        y' + p(x)y + q(x) = 0
    \]
    and introduce the IF:
    \[
        r(x) = \mye^{\int\myd{x}p(x)} \Rightarrow r'(x) = \mye^{\int\myd{x}p(x)}p(x) = r(x)p(x)
    \]
    multiplying the ODE by $r(x)$ gives
    \begin{align*}
        r(x)[y' + p(x)y + q(x)] &= 0 \\
        r(x)y' + \underbrace{r(x)p(x)}_{r'(x)}y + r(x)q(x) &= 0 \\
        [r(x)y]' + r(x)q(x) &= 0
    \end{align*}
    which can be solve by integration.
    % NB:  the IF should be regarded as a mathematical trick.  The integral is indefinite and can carry any arbitrary constant.  This additive constant becomes a multiplicative constant for $r(x)$.  We will see shortly that this changing this arbitrary constant will not ruin the math, but it will cancel in the final expression, so we are free to ignore it.

    % We currently get
    % \begin{align*}
    %     r(x)y &= -\int\myd{x}r(x)q(x) \\
    %     y &= -\frac{1}{r(x)}\int\myd{x}r(x)q(x)
    % \end{align*}
    \item for our system, we have
    \[
        y = p_0 \quad p(t) = k = \tau_\mathrm{rel}^{-1} \quad q(t) = -k_-
        \quad r(t) = \mye^{kt}
    \]
    so
    \begin{align*}
        \myder{}{t}\Big[\mye^{kt}p_0(t)\Big] &= k_- \mye^{kt} \\
        \Big[\mye^{kt'}p_0(t')\Big]_0^t &= k_- \int_0^t\myd{t'}\mye^{kt'} \\
        \mye^{kt}p_0(t) - p_0(0) &= \frac{k_-}{k}\Big[\mye^{kt} - 1\Big] \\
        p_0(t) &= \mye^{-kt}p_0(0) + \frac{k_-}{k}\Big[1 - \mye^{-kt}\Big]
    \end{align*}
    \item fluxes
    \[
        \mathrm{flux}_+ = k_+p_0 \quad \mathrm{flux}_- = k_-p_1
    \]
    \item detailed balance
    \[
        k_+p_0^\mathrm{eq} = k_-p_1^\mathrm{eq} \quad\Rightarrow\quad
        \frac{p_1^\mathrm{eq}}{p_0^\mathrm{eq}} = \frac{k_+}{k_-} = K_{eq}
    \]
    Notice that
    \[
        \frac{k_-}{k} = [1 + K_{eq}]^{-1}
    \]
    and
    \[
        \frac{p_1^\mathrm{eq}}{p_0^\mathrm{eq}}
        = \frac{1 - p_0^\mathrm{eq}}{p_0^\mathrm{eq}}
        = \frac{1}{p_0^\mathrm{eq}} - 1 = K_{eq}
    \]
    so
    \[
        p_0^\mathrm{eq} = [1 + K_{eq}]^{-1} = \frac{k_-}{k}
    \]
    Final expressions:  this is the familiar fluctuation-dissipation result for relaxation to equilibrium
    \begin{align*}
        p_0(t) - p_0^\mathrm{eq}
        &= \Big[p_0(0) - p_0^\mathrm{eq}\Big]\mye^{-kt} \\
        p_1(t) - p_1^\mathrm{eq}
        &= \Big[p_1(0) - p_1^\mathrm{eq}\Big]\mye^{-kt}
    \end{align*}
    \item critical result of nonequilibrium statistical mechanics:  systems relax to equilibrium exponentially.
    \item the timescale of relaxation is given by the transition rates:
    \[
        \Delta X\,\mye^{-kt} = \Delta X\,\mye^{-(k_+ + k_-)t}
        = \Delta X\,\mye^{-t/\tau_\mathrm{rel}}
    \]
\end{itemizec}

\subsection{Two-state system with transition matrix}
\begin{itemizec}
    \item rewrite time evolution wrt lag $\tau$ after some arbitrary time $t$
    \begin{align*}
        p_0(t+\tau)
        &= p_0^\mathrm{eq} + \Big[p_0(t) - p_0^\mathrm{eq}\Big]\mye^{-k\tau} \\
        p_1(t+\tau)
        &= p_1^\mathrm{eq} + \Big[p_1(t) - p_1^\mathrm{eq}\Big]\mye^{-k\tau}
    \end{align*}
    \item for each equation, multiply the second term by
    \[
        1 = p_0(t) + p_1(t)
    \]
    \item substitute the definitions of
    \[
        p_0^\mathrm{eq} \quad p_1^\mathrm{eq}
    \]
    in terms of the rate constants.
    \item some algebra will lead to an explicit expression for the transition matrix, $T_\tau$, such that
    \[
        p(t+\tau) =  T(\tau) p(t)
    \]
    with
    \[
        T(\tau) = k^{-1}\begin{bmatrix}
            k_- + k_+\mye^{-k\tau} & k_-(1-\mye^{-k\tau})  \\
            k_+(1-\mye^{-k\tau}) & k_+ + k_-\mye^{-k\tau}
        \end{bmatrix}
    \]
    being the \emph{left stochastic matrix} and $p(t)$ is the \emph{column vector} of probability.
\end{itemizec}

\subsection{Transition matrix for multistate systems}
\begin{itemizec}
    \item the equation is
    \[
        p_i(t+\tau) =  \sum_j T_{ij}(\tau) p_j(t)
    \]
    and $T$ is column-stochastic
    \[
        \sum_j T_{ij}(\tau) = 1
    \]
    with diagonal elements (``self-transitions'')
    \[
        T_{ii} = 1 - \sum_{j\ne i}T_{ji}
    \]
\end{itemizec}

\subsection{Short \& long time limits, two state}
\begin{itemizec}
    \item in the short time limit
    \[
        \lim_{\tau\rightarrow 0}T(\tau) =
        \begin{bmatrix}
            1 - k_+\tau & k_-\tau \\ k_+\tau & 1 - k_-\tau
        \end{bmatrix}
    \]
    \item in the long time limit, the transition matrix preserves equilibrium
    \[
        \lim_{\tau\rightarrow \infty}T(\tau) =
        \begin{bmatrix}
            k_-/k & k_-/k \\ k_+/k & k_+/k
        \end{bmatrix}
    \]
\end{itemizec}

\subsection{Markovian behaviour}
\begin{itemizec}
    \item if you can solve for the elements of a stochatic matrix
    \[
        T_{ij}(\tau) = \myP{x_{t+\tau}\in s_i | x_t\in s_j}
    \]
    \item the Chapman-Kolmogorov test for Markovianity provides a method for propagating the system to arbitrary multiples of $\tau$
    \[
        p(t+k\tau) =  T(\tau)^k p(t)
    \]
\end{itemizec}

\section{Probability and kinetics}

Taken from notes \texttt{07.22.cheat\_sheet}.
Want a place to put main kinetics result, including timescales, first passage times, etc.  Focus is on the two state model, but generalization where appropriate.

\subsection{The rate constant}
The fundamental definition of the rate constant is
\[
    k_{ij} = \lim_{\delta t\rightarrow 0}\frac{1}{\delta t}
    \myP{x_{t+\delta t} \in s_j|x_t\in s_i}
\]

% \appendix
% \input{sec-xxx}

\end{document}



% \usepackage{minted}
% \begin{minted}[mathescape, linenos]{python}
%     # Note: $\pi=\lim_{n\to\infty}\frac{P_n}{d}$
%     title = "Hello World"
% \end{minted}

%%=======================================================================
%%     PREFERRED FIGURE/TABLE
%%=======================================================================
%%                                 %%-------------------------------------%%
%% \begin{figure}[h!]              %%  caption right of figure
%%   \begin{minipage}[c]{0.67\textwidth}
%%     \includegraphics[width=\textwidth]{imgfile.png}
%%   \end{minipage}\hfill
%%   \begin{minipage}[c]{0.3\textwidth}
%%     \caption{Caption on the right hand side}
%%     \label{fig:}
%%   \end{minipage}
%% \end{figure}
%%                                 %%-------------------------------------%%
%% \begin{figure}[h!]\centering    %%  caption below figure
%%   \includegraphics[width=0.5\textwidth]{imgfile.png}
%%   \caption{General caption}
%%   \label{fig:label}
%% \end{figure}
%%                                 %%-------------------------------------%%
%% \begin{figure}[!h]\centering    %%  array of figures with subcaptions
%%   \subfloat[]{\includegraphics[width=0.5\textwidth]{}}
%%   \subfloat[]{\includegraphics[width=0.5\textwidth]{}} \\
%%   \subfloat[]{\includegraphics[width=0.5\textwidth]{}}
%%   \subfloat[]{\includegraphics[width=0.5\textwidth]{}} \\
%%   \subfloat[]{\includegraphics[width=0.5\textwidth]{}}
%%   \subfloat[]{\includegraphics[width=0.5\textwidth]{}}
%%   \caption{General caption}
%%   \label{fig:label}
%% \end{figure}
%%                                        %%------------------------------%%
%% \begin{table}[!h]\centering\ra{1.3}    %%  simple table
%%   \begin{tabular}{@{}rrrr@{}}\toprule
%%     & $t=0$ & $t=1$ & $t=2$ \\ \midrule
%%     $c$ & 1 & 2 & 3 \\
%%     $c$ & 1 & 2 & 3 \\
%%     \bottomrule
%%   \end{tabular}
%%   \caption{Caption}
%% \end{table}
%%                                        %%------------------------------%%
%% \begin{table}[!h]\centering\ra{1.3}    %%  complicated table
%%   \begin{tabular}{@{}rrrrcrrr@{}}\toprule
%%     & \multicolumn{3}{c}{$w = 8$} & \phantom{abc}&
%%     \multicolumn{3}{c}{$w = 16$}\\
%%     \cmidrule{2-4} \cmidrule{6-8} \cmidrule{10-12}
%%     & $t=0$ & $t=1$ & $t=2$ && $t=0$ & $t=1$ & $t=2$ \\ \midrule
%%     \rowcolor[gray]{.95} $c$ & 1 & 2 & 3 && 4 & 5 & 6 \\
%%     \rowcolor[gray]{1.0} $c$ & 1 & 2 & 3 && 4 & 5 & 6 \\
%%     \bottomrule
%%   \end{tabular}
%%   \caption{Caption}
%% \end{table}

%%=======================================================================
%%     OLDER FIGURE/TABLE
%%=======================================================================
%% \begin{center}
%%   \includegraphics[width=0.75\linewidth]{}
%%   \captionof{figure}{}
%% \end{center}
%%
%% \begin{center}
%%   \begin{tabular}{rl}
%%     \includegraphics[width=0.5\linewidth]{} &
%%     \includegraphics[width=0.5\linewidth]{} \\
%%     \includegraphics[width=0.5\linewidth]{} &
%%     \includegraphics[width=0.5\linewidth]{} \\
%%   \end{tabular}
%%   \captionof{figure}{}
%% \end{center}
%%
%% \begin{center}
%%   \begin{tabular}{lll}
%%     \hline
%%     A & B & C \\
%%     \hline
%%     \hline
%%     1 & 2 & 3 \\
%%     \hline
%%   \end{tabular}
%%   \captionof{table}{test}
%% \end{center}
